\chapter{Bilirubin (Urine)}

\section*{What Is This Test?}
A urinary bilirubin test detects conjugated bilirubin that has escaped from the liver into the bloodstream and is filtered by the kidneys. Unconjugated bilirubin is bound to albumin and cannot be filtered by the kidney, so it never appears in urine. The presence of urinary bilirubin therefore confirms that conjugated bilirubin has reached the systemic circulation -- indicating hepatic or post-hepatic pathology (hepatocellular or obstructive jaundice). The absence of urinary bilirubin in a jaundiced patient suggests unconjugated hyperbilirubinemia (hemolysis, Gilbert syndrome). The test is most commonly performed as part of the routine urinalysis dipstick.

\section*{Alternative Names}
Urine Bilirubin; Urinalysis Bilirubin; Bilirubinuria; Dipstick Bilirubin

\section*{Principle}
The diazo reaction on the urinary dipstick pad couples bilirubin with 2,4-dichloroaniline or 2,6-dichlorobenzenediazonium tetrafluoroborate to form a colored azo pigment. Color intensity is read visually or by automated urinalysis readers and graded as negative, small (+), moderate (++), or large (+++). Detection threshold is approximately 0.4 mg/dL. False negatives may occur with vitamin C and with prolonged exposure to light.

\section*{History}
Urinary bilirubin was one of the earliest recognized diagnostic signs of jaundice, documented by van den Bergh in the 1910s \cite{vandenbergh1916uber}. The diazo-based dipstick pad was introduced by Helen and Alfred Free in the 1950s as part of the multitest dipstick.

\section*{Why This Test Is Ordered}
\begin{itemize}
  \item Evaluation of jaundice and suspected biliary obstruction
  \item Differentiating conjugated from unconjugated hyperbilirubinemia
  \item Monitoring hepatic dysfunction (hepatitis, drug toxicity, cirrhosis)
  \item Routine screening during urinalysis
  \item Evaluation of suspected cholelithiasis or choledocholithiasis
\end{itemize}

\section*{Primary vs Secondary Test}
Primary screening test; results guide serum bilirubin (total and direct) fractionation and liver imaging.

\section*{How This Test Is Performed}
A fresh random urine sample is collected; midstream clean-catch is preferred. The dipstick pad is briefly immersed, and the color is read at the time specified (typically 30--60 seconds). The specimen must be fresh because bilirubin degrades to biliverdin on exposure to light and oxidation. Confirmatory testing is done by urine microscopy or via serum testing.

\section*{What Preparation Is Needed}
None. Hydration is normal. Patient should not take large-dose vitamin C prior.

\section*{Contraindications and Precautions}
False negatives: photo-oxidation (specimen sitting in light $>$ 1 hour), high-dose vitamin C, high urinary nitrite. False positives: phenazopyridine (Pyridium), indole, large doses of chlorpromazine metabolites, rifampin (also discolors urine). Specimen should be protected from light.

\section*{Risks}
None.

\section*{What Happens After the Test}
Dipstick results are immediately available. If bilirubin is positive, the laboratory reflexively performs microscopic examination and the clinician orders serum total/direct bilirubin and liver function tests.

\section*{Understanding Results}

\subsection*{Below Normal}
Negative urinary bilirubin is normal. If the patient is clinically jaundiced but urine bilirubin is absent, suspect unconjugated hyperbilirubinemia (hemolysis, Gilbert, Crigler-Najjar). If the jaundice is conjugated, the absence may be due to specimen light exposure or a false negative.

\subsection*{Above Normal}
\begin{itemize}
  \item Obstructive jaundice: common bile duct stone, stricture, pancreatic or biliary malignancy, cholangiocarcinoma
  \item Hepatocellular disease: viral hepatitis, drug-induced liver injury, cirrhosis, alcoholic hepatitis
  \item Sepsis and hepatic failure
  \item Bilirubinuria often appears before clinical jaundice is visible, making it an early marker
  \item Conjugated bilirubin also causes dark-colored urine (think ``Coca-Cola urine'')
\end{itemize}

\section*{Normal Results}
\begin{itemize}
  \item \textbf{Bilirubin (urine, dipstick):} Negative
  \item \textbf{If jaundiced patient:} Negative urine bilirubin suggests unconjugated (pre-hepatic) hyperbilirubinemia
  \item \textbf{If positive:} urine bilirubin suggests conjugated (hepatic or post-hepatic) hyperbilirubinemia
\end{itemize}

\section*{Follow-up Tests}
\begin{itemize}
  \item \textbf{Serum bilirubin (total and direct):} Confirms hyperbilirubinemia and fractionation
  \item \textbf{Urine urobilinogen:} Classifies jaundice (elevated in hemolysis and hepatocellular; absent in biliary obstruction)
  \item \textbf{Liver function tests (ALT, AST, ALP, GGT):} Pattern differentiates hepatic vs biliary disease
  \item \textbf{Abdominal ultrasound:} First-line imaging for biliary obstruction
  \item \textbf{MRCP or ERCP:} Definitive imaging and therapeutic intervention for choledocholithiasis
  \item \textbf{Hepatitis serology:} If viral hepatitis suspected
\end{itemize}

\section*{References}
\bibliographystyle{unsrtnat}
\bibliography{references}

\clearpage
\section*{Regulatory Status and Authorization Date}
Regulatory status applies to the specific assay, instrument, manufacturer, and intended use rather than to the test name alone. No single approval date can be assigned to this test category. For a commercial assay, verify the current product in the relevant regulator's database: the U.S. Food and Drug Administration (FDA) clearance, approval, or authorization; India's Central Drugs Standard Control Organisation (CDSCO) licence or permission; the European Union In Vitro Diagnostic Regulation (EU IVDR) CE certificate issued through the applicable conformity-assessment route; or the United Kingdom Medicines and Healthcare products Regulatory Agency (MHRA) registration and UKCA/CE status. Laboratory-developed tests may not have an individual FDA, CDSCO, EU IVDR, or MHRA approval date.
\clearpage
